\documentclass[12pt]{article}
\usepackage[T1]{fontenc}
\usepackage[utf8]{inputenc}
\usepackage{lmodern}
\usepackage{textcomp}
\usepackage{enumitem}
\usepackage{geometry}
\geometry{margin=1in}
\begin{document}
\begin{center}
Answer Key - Political Science - Undergraduate Level Assessment 19 \
\small (Answers are ordered exactly as in the question paper.)
\end{center}

\section*{Multiple Choice}
\begin{enumerate}[label=\arabic*.]
\item \textbf{Answer:} C
\\ \textit{Options:} A) terrorism policy.; B) economic policy.; C) foreign policy.; D) international policy.
\\ \textit{Note:} The correct answer is "foreign policy" because it specifically refers to the strategies and decisions made by a nation regarding its interactions and relationships with other countries.
\item \textbf{Answer:} D
\\ \textit{Options:} A) negotiate treaties; B) veto legislation; C) grant reprieves and pardons; D) declare war
\\ \textit{Note:} The power to declare war is not constitutionally mandated for the president, as the U.S. Constitution grants this authority specifically to Congress, while the president serves as the Commander-in-Chief of the armed forces.
\item \textbf{Answer:} B
\\ \textit{Options:} A) The electoral college; B) Political parties; C) Separation of powers; D) The term length for members of Congress
\\ \textit{Note:} The Constitution does not mention political parties, reflecting the Founding Fathers' intention to avoid factionalism in the new government.
\item \textbf{Answer:} D
\\ \textit{Options:} A) Foreign policies of other states; B) International law; C) Intergovernmental organizations; D) All of the above
\\ \textit{Note:} The correct answer is "All of the above" because US foreign policy decision-making is influenced by a variety of constraints, including domestic political considerations, international relations, economic factors, and public opinion, all of which can significantly impact policy outcomes.
\item \textbf{Answer:} D
\\ \textit{Options:} A) It ended the influence of American exceptionalism entirely; B) Exceptionalism was enhanced by America's status as the 'leader of the free world'; C) The extension of American power globally challenged core assumptions of exceptionalism; D) Both b and c
\\ \textit{Note:} The implications of the Cold War for American exceptionalism included the reinforcement of the belief in the United States as a global leader in promoting democracy and capitalism, while also highlighting the ideological struggle against communism, which collectively shaped national identity and foreign policy.
\end{enumerate}

\section*{True / False}
\begin{enumerate}[label=\arabic*.]
\setcounter{enumi}{5}
\item \textbf{Answer:} True
\\ \textit{Options:} A) True; B) False
\\ \textit{Note:} The term 'iron triangle' is true because it describes the enduring and mutually beneficial relationships formed between federal agencies, congressional committees, and interest groups or lobbyists, which work together to influence policy and secure favorable outcomes for their respective interests.
\item \textbf{Answer:} False
\\ \textit{Options:} A) True; B) False
\\ \textit{Note:} The statement is false because the Republican Party has traditionally opposed labor unions and has sought to reduce their influence in politics, particularly since the 1980 s, when it shifted toward a more conservative and pro-business agenda.
\item \textbf{Answer:} False
\\ \textit{Options:} A) True; B) False
\\ \textit{Note:} The Voting Rights Act of 1965 prohibited literacy tests as a prerequisite for voting, aiming to eliminate discriminatory practices that disenfranchised voters, particularly African Americans.
\item \textbf{Answer:} False
\\ \textit{Options:} A) True; B) False
\\ \textit{Note:} Strict constructionism maintains that the Bill of Rights is designed to protect individual liberties against both federal and state encroachments, thus applying its limitations to both levels of government.
\item \textbf{Answer:} True
\\ \textit{Options:} A) True; B) False
\\ \textit{Note:} The Voting Rights Act of 1965 prohibited literacy tests as a means to eliminate discriminatory practices that hindered voter registration and access, particularly for African Americans in the South.
\end{enumerate}

\section*{Long Form Response}
\begin{enumerate}[label=\arabic*.]
\setcounter{enumi}{10}
\item \textbf{Answer:} In The Federalist No. 10, James Madison argues that the existence of factions, or groups of citizens united by a common interest, poses a threat to the stability and fairness of government. He acknowledges that while factions are inevitable due to the diverse nature of human interests, their negative effects can be mitigated through a republican form of government. Madison contends that a large republic, with a diverse array of competing interests, will dilute the influence of any single faction, thus preventing any one group from dominating the political landscape. Furthermore, by electing representatives who are informed and accountable to a wider constituency, a republican government can ensure that policies reflect the common good rather than the interests of a specific faction. In this way, Madison highlights the importance of a structured political system to manage factionalism and promote a balanced governance.
\\ \textit{Note:} correct because it accurately captures Madison's central argument that while factions are a natural consequence of varying human interests, a large republican government can effectively reduce the dominance of any single faction by promoting a multiplicity of competing interests, thereby enhancing stability and fairness in governance.
\item \textbf{Answer:} The geopolitical dynamics of the Cold War significantly shaped U.S. perceptions and policies towards developments in the Third World, as local events were often interpreted through the lens of superpower rivalry. For instance, the U.S. supported anti-communist regimes and movements, such as in Vietnam and Latin America, viewing the spread of communism as a direct threat to American interests. This led to interventions, such as the Bay of Pigs invasion in Cuba and support for authoritarian regimes in countries like Chile, under the belief that stabilizing these regions would prevent the expansion of Soviet influence. Consequently, U.S. foreign policy prioritized ideological alignment over democratic values, demonstrating how Cold War tensions redefined its approach to global affairs.
\\ \textit{Note:} correct because it effectively illustrates how the Cold War's superpower rivalry drove U.S. foreign policy in the Third World, emphasizing the perception of communism as a threat and providing specific examples of U.S. interventions and support for anti-communist regimes.
\end{enumerate}

\vfill
\noindent\textit{End of Answer Key}
\end{document}